% ============================================================================
% BIBLIOGRAFIA
% ============================================================================
% Template SGGW używa ręcznego \thebibliography (nie BibTeX/biblatex).
% Wpisy uporządkowane tematycznie: biomechanika biegu / pose estimation / ML / metodologia.
%
% Status: 27 wpisów (sesja 2026-05-23). Zweryfikowane 3x (agenci) + ręcznie (user).
% Diaz 2019 usunięty — niepotwierdzony. OpenPose rok zmieniony na 2021 (tom 43, nr 1).
% Materiał źródłowy: docs/thesis-notes/2026-05-14-research-podobnych-prac.md
%                    docs/reference-values.md (cytaty biomechaniczne w rules.py)

\begin{thebibliography}{99}

% ============================================================================
% Pose estimation z monocular 2D wideo
% ============================================================================

\bibitem{bazarevsky2020blazepose}
V.\ Bazarevsky, I.\ Grishchenko, K.\ Raveendran, T.\ Zhu, F.\ Zhang, M.\ Grundmann,
\textit{BlazePose: On-device Real-time Body Pose Tracking},
CVPR Workshop on Computer Vision for Augmented and Virtual Reality, arXiv:2006.10204, 2020.

\bibitem{cao2021openpose}
Z.\ Cao, G.\ Hidalgo, T.\ Simon, S.E.\ Wei, Y.\ Sheikh,
\textit{OpenPose: Realtime Multi-Person 2D Pose Estimation Using Part Affinity Fields},
IEEE Transactions on Pattern Analysis and Machine Intelligence, 43(1): 172--186, 2021.
\url{https://doi.org/10.1109/TPAMI.2019.2929257}

\bibitem{lin2014coco}
T.Y.\ Lin et al.,
\textit{Microsoft COCO: Common Objects in Context},
European Conference on Computer Vision (ECCV), pp.\ 740--755, 2014.
\url{https://doi.org/10.1007/978-3-319-10602-1_48}

% ============================================================================
% Walidacja 2D pose vs 3D motion capture (gait/running)
% ============================================================================

\bibitem{stenum2021}
J.\ Stenum, C.\ Rossi, R.T.\ Roemmich,
\textit{Two-dimensional video-based analysis of human gait using pose estimation},
PLOS Computational Biology, 17(4): e1008935, 2021.
\url{https://doi.org/10.1371/journal.pcbi.1008935}

\bibitem{ali2024dmd}
M.M.\ Ali, M.M.\ Hassan, M.\ Zaki,
\textit{Human Pose Estimation for Clinical Analysis of Gait Pathologies},
Bioinformatics and Biology Insights, 2024.
\url{https://doi.org/10.1177/11779322241231108}

\bibitem{menychtas2023}
D.\ Menychtas, N.\ Petrou, I.\ Kansizoglou, E.\ Giannakou, A.\ Grekidis, A.\ Gasteratos,
V.\ Gourgoulis, E.\ Douda, I.\ Smilios, M.\ Michalopoulou, G.C.\ Sirakoulis, N.\ Aggeloulis,
\textit{Gait analysis comparison between manual marking, 2D pose estimation algorithms, and 3D marker-based system},
Frontiers in Rehabilitation Sciences, 4: 1238134, 2023.
\url{https://doi.org/10.3389/fresc.2023.1238134}

\bibitem{hgcnmlp2024}
R.\ Hu, Y.\ Diao, Y.\ Wang, G.\ Li, R.\ He, Y.\ Ning, N.\ Lou, G.\ Li, G.\ Zhao,
\textit{Effective evaluation of HGcnMLP method for markerless 3D pose estimation of musculoskeletal diseases patients based on smartphone monocular video},
Frontiers in Bioengineering and Biotechnology, 11: 1335251, 2024.
% Uwaga: volume 11 (2023), ale publikacja online 9 stycznia 2024
\url{https://doi.org/10.3389/fbioe.2023.1335251}

\bibitem{frontiers2025review}
S.\ Edriss, C.\ Romagnoli, L.\ Caprioli, V.\ Bonaiuto, E.\ Padua, G.\ Annino,
\textit{Commercial vision sensors and AI-based pose estimation frameworks for markerless motion analysis in sports and exercises: a mini review},
Frontiers in Physiology, 16: 1649330, 2025.
\url{https://doi.org/10.3389/fphys.2025.1649330}

\bibitem{roggio2024review}
F.\ Roggio, B.\ Trovato, M.\ Sortino, G.\ Musumeci,
\textit{A comprehensive analysis of the machine learning pose estimation models used in human movement and posture analyses: A narrative review},
Heliyon, 10(21): e39977, 2024.
\url{https://doi.org/10.1016/j.heliyon.2024.e39977}

\bibitem{vincent2024footstrike}
H.K.\ Vincent, K.\ Coffey, A.\ Villasuso, K.R.\ Vincent, S.\ Sharififar, L.\ Pezzullo, R.M.\ Nixon,
\textit{Accuracy of self-reported foot strike pattern detection among endurance runners},
Frontiers in Sports and Active Living, 6: 1491486, 2024.
\url{https://doi.org/10.3389/fspor.2024.1491486}

% ============================================================================
% Biomechanika biegu — reference values dla reguł rekomendacji
% ============================================================================

\bibitem{heiderscheit2011}
B.C.\ Heiderscheit, E.S.\ Chumanov, M.P.\ Michalski, C.M.\ Wille, M.B.\ Ryan,
\textit{Effects of step rate manipulation on joint mechanics during running},
Medicine \& Science in Sports \& Exercise, 43(2): 296--302, 2011.
\url{https://doi.org/10.1249/MSS.0b013e3181ebedf4}

\bibitem{novacheck1998}
T.F.\ Novacheck,
\textit{The biomechanics of running},
Gait \& Posture, 7(1): 77--95, 1998.
\url{https://doi.org/10.1016/S0966-6362(97)00038-6}

\bibitem{souza2016}
R.B.\ Souza,
\textit{An Evidence-Based Videotaped Running Biomechanics Analysis},
Physical Medicine and Rehabilitation Clinics of North America, 27(1): 217--236, 2016.
\url{https://doi.org/10.1016/j.pmr.2015.08.006}

% Diaz 2019 (vertical oscillation) — usunięty: nie udało się potwierdzić istnienia pracy.
% Jeśli znajdziesz pełne dane bibliograficzne, dodaj wpis z powrotem.

\bibitem{robinson1987}
R.O.\ Robinson, W.\ Herzog, B.M.\ Nigg,
\textit{Use of force platform variables to quantify the effects of chiropractic manipulation on gait symmetry},
Journal of Manipulative and Physiological Therapeutics, 10(4): 172--176, 1987.

\bibitem{daoud2012}
A.I.\ Daoud et al.,
\textit{Foot strike and injury rates in endurance runners: a retrospective study},
Medicine \& Science in Sports \& Exercise, 44(7): 1325--1334, 2012.
\url{https://doi.org/10.1249/MSS.0b013e3182465115}

% ============================================================================
% Uczenie maszynowe / klasyfikatory faz biegu
% ============================================================================

\bibitem{jeon2024bilstm}
H.\ Jeon, D.\ Lee,
\textit{Bi-Directional Long Short-Term Memory-Based Gait Phase Recognition Method Robust to Directional Variations in Subject's Gait Progression Using Wearable Inertial Sensor},
Sensors, 24(4): 1276, 2024.
\url{https://doi.org/10.3390/s24041276}

\bibitem{su2020dcnn}
B.\ Su, C.\ Smith, E.\ Gutierrez~Farewik,
\textit{Gait Phase Recognition Using Deep Convolutional Neural Network with Inertial Measurement Units},
Biosensors, 10(9): 109, 2020.
\url{https://doi.org/10.3390/bios10090109}

\bibitem{vu2020review}
H.T.T.\ Vu, D.\ Dong, H.-L.\ Cao, T.\ Verstraten, D.\ Lefeber, B.\ Vanderborght, J.\ Geeroms,
\textit{A Review of Gait Phase Detection Algorithms for Lower Limb Prostheses},
Sensors, 20(14): 3972, 2020.
\url{https://doi.org/10.3390/s20143972}

\bibitem{dimmick2023fatigue}
H.L.\ Dimmick, C.R.\ van~Rassel, M.J.\ MacInnis, R.\ Ferber,
\textit{Use of subject-specific models to detect fatigue-related changes in running biomechanics: a random forest approach},
Frontiers in Sports and Active Living, 5: 1283316, 2023.
\url{https://doi.org/10.3389/fspor.2023.1283316}

\bibitem{martinez2020rf}
J.\ Mart\'{i}nez-Gramage, J.P.\ Albiach, I.N.\ Molt\'{o}, J.J.\ Amer-Cuenca, V.H.\ Moreno, E.\ Segura-Ort\'{i},
\textit{A Random Forest Machine Learning Framework to Reduce Running Injuries in Young Triathletes},
Sensors, 20(21): 6388, 2020.
\url{https://doi.org/10.3390/s20216388}

\bibitem{young2023smartphone}
F.\ Young, R.\ Mason, R.\ Morris, S.\ Stuart, A.\ Godfrey,
\textit{Internet-of-Things-Enabled Markerless Running Gait Assessment from a Single Smartphone Camera},
Sensors, 23(2): 696, 2023.
\url{https://doi.org/10.3390/s23020696}

% ============================================================================
% Filtrowanie sygnałów biomechanicznych
% ============================================================================

\bibitem{crenna2021filtering}
F.\ Crenna, G.B.\ Rossi, M.\ Berardengo,
\textit{Filtering Biomechanical Signals in Movement Analysis},
Sensors, 21(13): 4580, 2021.
\url{https://doi.org/10.3390/s21134580}

\bibitem{savitzky1964}
A.\ Savitzky, M.J.E.\ Golay,
\textit{Smoothing and Differentiation of Data by Simplified Least Squares Procedures},
Analytical Chemistry, 36(8): 1627--1639, 1964.
\url{https://doi.org/10.1021/ac60214a047}

% ============================================================================
% Prewencja urazów biegowych / gait retraining
% ============================================================================

\bibitem{figueiredo2025cadence}
I.\ Figueiredo, M.\ Reis~e~Silva, J.E.\ Sousa,
\textit{The Influence of Running Cadence on Biomechanics and Injury Prevention: A Systematic Review},
Cureus, 17(8): e90322, 2025.
\url{https://doi.org/10.7759/cureus.90322}

% ============================================================================
% Uczenie maszynowe — podręczniki
% ============================================================================

\bibitem{goodfellow2016deep}
I.\ Goodfellow, Y.\ Bengio, A.\ Courville,
\textit{Deep Learning},
MIT Press, 2016.

\bibitem{hochreiter1997lstm}
S.\ Hochreiter, J.\ Schmidhuber,
\textit{Long Short-Term Memory},
Neural Computation, 9(8): 1735--1780, 1997.
\url{https://doi.org/10.1162/neco.1997.9.8.1735}

\bibitem{breiman2001rf}
L.\ Breiman,
\textit{Random Forests},
Machine Learning, 45(1): 5--32, 2001.
\url{https://doi.org/10.1023/A:1010933404324}

% ============================================================================
% Narzędzia / framework
% ============================================================================

\bibitem{pytorch}
A.\ Paszke et al.,
\textit{PyTorch: An Imperative Style, High-Performance Deep Learning Library},
Advances in Neural Information Processing Systems (NeurIPS), 2019.

\bibitem{scipy}
P.\ Virtanen et al.,
\textit{SciPy 1.0: fundamental algorithms for scientific computing in Python},
Nature Methods, 17: 261--272, 2020.
\url{https://doi.org/10.1038/s41592-019-0686-2}

\bibitem{sklearn}
F.\ Pedregosa et al.,
\textit{Scikit-learn: Machine Learning in Python},
Journal of Machine Learning Research, 12: 2825--2830, 2011.

\bibitem{opencv}
G.\ Bradski,
\textit{The OpenCV Library},
Dr. Dobb's Journal of Software Tools, 2000.

\end{thebibliography}
